% Generated draft: requires mathematical and source review.
\documentclass[11pt]{article}
\usepackage[margin=25mm]{geometry}
\usepackage{fontspec}
\setmainfont{Latin Modern Roman}
\usepackage{amsmath,amssymb,mathtools,unicode-math}
\setmathfont{Latin Modern Math}
\usepackage{enumitem,needspace,xurl,hyperref}
\setlength{\parindent}{0pt}
\setlength{\parskip}{5pt}
\setlength{\emergencystretch}{3em}
\newcommand{\N}{\mathbb{N}}
\newcommand{\Z}{\mathbb{Z}}
\newcommand{\Q}{\mathbb{Q}}
\newcommand{\R}{\mathbb{R}}
\newcommand{\C}{\mathbb{C}}
\newcommand{\F}{\mathbb{F}}
\newcommand{\A}{\mathbb{A}}
\newcommand{\PP}{\mathbb{P}}
\newcommand{\E}{\mathbb{E}}
\newcommand{\eps}{\varepsilon}
\newcommand{\dd}{\,\mathrm{d}}
\begin{document}
\section*{UnsolvedMath ID: 30002932}
\textbf{Source problem code:} catalogue entry
\subsection{Definitions and mathematical statement}
For each integer \(n\geq 1\), let \(S_n\) denote the permutations of \(\{1,\ldots,n\}\). Write \(\sigma_n\in S_n\) for the permutation at stage \(n\) of the Chinese restaurant process growth referred to in the supplied record. Define the number of fixed points by
\[
A_1(n)=\#\{i\in\{1,\ldots,n\}:\sigma_n(i)=i\}.
\]
For a realization of this growth process, define its zero set and zero-occupation count by
\[
\mathcal Z=\{n\geq1:A_1(n)=0\},\qquad
Z_N=\sum_{n=1}^{N}\mathbf 1_{\{A_1(n)=0\}},\quad N\geq1.
\]
The supplied motivating question asks for pathwise properties of \((A_1(n))_{n\geq1}\), particularly how often it vanishes. Possible precise interpretations include determining whether \(\mathcal Z\) is infinite almost surely, or determining the almost-sure asymptotic behavior of \(Z_N\) and the existence and value of \(\lim_{N\to\infty}Z_N/N\). These are interpretations of the wording, not separately source-stated conjectures.

\textbf{Unresolved formulation qualifications.} The supplied statement does not specify the Chinese restaurant process parameter, transition probabilities, or initialization, and does not select a meaning of ``how often'' or a precise collection of pathwise properties. Accordingly, no particular probability law, limiting value, or universal almost-sure assertion is imposed here. The supplied verification marks the statement as \emph{unrecoverable}: it reports that the official report uses this question as motivation and subsequently gives relevant time-change limit results, without isolating a separate unresolved claim. This draft therefore records the motivating question and its definable quantities; it does not promote it to a verified open conjecture.
\subsection{Short English statement}
Track the number of fixed points along the permutation growth process and ask how often that number is zero. The process law and the intended notion of frequency require clarification. The record labels the topic partially solved, but its supplied source verification identifies a motivating question rather than a separately stated open problem.
\subsection{Sources}
\begin{itemize}
\item Local registry: \texttt{data/unsolvedmath/problems.json}, permanent UnsolvedMath ID \(30002932\), source code \texttt{catalogue entry}; title: \emph{Motivating Fixed-Point Question for Chinese Restaurant Permutations}. The supplied record gives proposed year 2015 and names no individual proposer.
\item Supplied report citation: \emph{Interplay of Analysis and Probability in Applied Mathematics (2016)}, Oberwolfach Reports, pp.\ 1989--2064, DOI \texttt{10.4171/owr/2015/35}. \url{https://doi.org/10.4171/owr/2015/35}. The record describes Oberwolfach Reports as documenting mathematical workshops and their open-problem discussions. The citation uses 2016, whereas the DOI contains 2015 and the proposed-year field is 2015; this draft preserves those supplied details without resolving their relationship.
\item Supplied literature reference: Gnedin--Stark, \emph{Random permutations and queues}. \url{https://arxiv.org/abs/2209.12781}.
\item Supplied dated literature curation, checked \texttt{2026-08-21}: status \emph{partially solved}. Its assessment says that Gnedin and Stark resolve the principal time-average questions for fixed points in the discrete Chinese restaurant process and its continuous-time embedding, while the broad request for all pathwise properties of \(A_1\) remains open. This is supplied dated curation derived from the OWR material, not an independently verified literature assessment; its time-sensitive status and cited claims require independent verification.
\item Supplied statement-verification note, retained in the record updated \texttt{2026-09-16T12:06:53Z}: the official report is described as introducing this pathwise question as motivation and then giving relevant time-change limit results, rather than presenting a separate unresolved problem. The extraction does not isolate a source-stated open claim; no clean statement is promoted, and the original record remains retained and publicly visible for provenance. These are supplied verification claims, not claims of source inspection in preparing this draft.
\end{itemize}
\end{document}
